\chapter{Project Context and Scope}
\section*{Introduction}
This chapter presents the host organization, the internship framework, the
existing training context and the scope selected for the project. It explains
why Safyron Learning is relevant to the activities of Safyron Engineering and
how the one-month work was organized.

\section{Host Organization and Internship Framework}
The professional internship took place from 1 to 31 August 2026 at
\textbf{Safyron Engineering}, located in Rue Lac Toba, Les Berges du Lac,
Tunis. The company operates in the field of industrial technical studies,
including fire protection, hazard studies and industrial risk analysis. These
activities require accurate technical documentation and continuous professional
training.

\begin{figure}[H]\centering
\includegraphics[width=0.36\textwidth]{../elearning-frontend/public/Logo.png}
\caption{Safyron Engineering visual identity}
\label{fig:safyron-logo}
\end{figure}

The internship was completed as part of the fifth year of the software
engineering program at EPI Digital School Sousse. The work was carried out
individually and supervised professionally by Mr Ammar Achraf, founder of
Safyron Engineering. The company context made it possible to connect software
engineering practices with the operational need for reliable safety and
compliance training.

\section{Project overview}
\projectname{} is a compliance-oriented learning management platform developed
as a full-stack web application. Its visual identity is based on Safyron
Engineering and its functional vocabulary reflects industrial training:
mandatory modules, published resources, assessment attempts, audit records and
instructor approval.

The project combines a responsive Angular single-page application with a
Spring Boot REST backend. MySQL stores users, modules, materials, quizzes,
questions, learner progress and quiz results. Uploaded PDF, DOCX and PPTX
documents are retained as learning resources and can also provide the source
text for quiz generation.

\section{Initial problem}
Safyron Engineering's technical activities depend on documents, procedures and
regulatory knowledge that must remain available to employees. Before the
platform, training material, quiz preparation and evidence of completion could
be handled as separate activities. This separation made it difficult to know
which version of a module was current, whether a learner had completed the
assessment and whether the result had been reviewed by a responsible
instructor.

The project therefore models the training process as a complete lifecycle:
material is attached to a module, an assessment is prepared, the module is
published, the learner submits answers once, and the instructor records the
final decision. The administrator receives a transverse view for governance
and data maintenance.

Traditional training workflows often separate the content repository,
questionnaire authoring and evidence of completion. This creates several
problems:
\begin{itemize}
\item learners cannot easily identify the current published material;
\item instructors spend significant time rewriting questions from documents;
\item a score alone does not provide enough evidence for approval;
\item repeated attempts can invalidate the meaning of an assessment;
\item deleting a module or user can fail because of relational dependencies.
\end{itemize}

The project question was therefore: \emph{how can a web platform connect
learning resources, AI-assisted assessment and instructor validation while
preserving data consistency?}

\section{Stakeholders}
\subsection{Learner}
The learner browses published modules, opens or downloads files, takes an
available quiz and sees the submission state. The score is not exposed as a
learner-facing decision; certification is controlled by the instructor.

\subsection{Instructor}
The instructor creates and manages modules, imports source material, produces
AI-generated or manual quizzes, publishes assessments and reviews detailed
answers. Approval can be removed if it was granted accidentally.

\subsection{Administrator}
The administrator manages users, modules, quizzes and operational statistics.
The dashboard groups submissions by module and provides a global governance
view.

\section{Scope and objectives}
\begin{table}[H]\centering
\caption{Project objectives}\label{tab:objectives}
\begin{tabularx}{\textwidth}{>{\bfseries}p{0.28\textwidth}X}
\toprule Objective & Expected result\\\midrule
Module management & Create, publish, update and delete modules with dependencies.\\
Learning resources & Upload, preview, open and download training files.\\
Assessment & Generate quizzes with AI or create them manually.\\
Controlled attempt & Store one submission per learner and quiz.\\
Human validation & Display detailed answers and support approval or reversal.\\
Administration & Manage users, modules and grouped submission indicators.\\
Quality & Validate the API through tests, builds and Swagger.\\
\bottomrule
\end{tabularx}
\end{table}

\section{Planning}
The month was divided into four iterations. The first week focused on
requirements, architecture and the existing codebase. The second covered
modules, files and learner navigation. The third implemented quiz generation,
manual authoring, one-time submission and review. The fourth concentrated on
administration, deletion consistency, Swagger, interface improvements and
end-to-end validation.

\begin{table}[H]
\centering
\caption{Simplified Gantt diagram -- provisional internship planning}
\label{tab:gantt-planning}
\renewcommand{\arraystretch}{1.35}
\begin{tabular}{@{}p{0.58\textwidth}*{4}{>{\centering\arraybackslash}p{0.085\textwidth}}@{}}
\toprule
\textbf{Proposed activity}
  & \textbf{\shortstack{S1\\[-2pt]\footnotesize 01--07}}
  & \textbf{\shortstack{S2\\[-2pt]\footnotesize 08--14}}
  & \textbf{\shortstack{S3\\[-2pt]\footnotesize 15--21}}
  & \textbf{\shortstack{S4\\[-2pt]\footnotesize 22--31}}\\
\midrule
Understanding the organisational context
  & \rule{0.34cm}{0.34cm} & & & \\
Requirements analysis and scope definition
  & \rule{0.34cm}{0.34cm} & & & \\
Functional analysis and UML modelling
  & & \rule{0.34cm}{0.34cm} & & \\
Architecture and data-model design
  & & \rule{0.34cm}{0.34cm} & & \\
Backend implementation
  & & & \rule{0.34cm}{0.34cm} & \\
Progressive integration and finalisation
  & & & \rule{0.34cm}{0.34cm} & \rule{0.34cm}{0.34cm} \\
Testing and validation
  & & & & \rule{0.34cm}{0.34cm} \\
AI integration and final corrections
  & & & & \rule{0.34cm}{0.34cm} \\
\bottomrule
\end{tabular}
\end{table}

\section{Project boundaries}
The project does not attempt to provide payments, video conferencing,
multi-tenant billing or a complete certificate document generator. These
features are suitable future extensions. The internship focused on a reliable
learning and assessment core.

\section*{Conclusion}
The company context justifies a platform focused on traceable learning and
controlled assessment rather than a generic document repository. The next
chapter translates this context into functional and non-functional
requirements, architecture decisions and a persistent data model.
